\documentclass[12pt,a4paper]{article}

\usepackage{fontspec}
\usepackage{polyglossia}
\setmainlanguage{vietnamese}
\setmainfont{DejaVu Serif}
\setsansfont{DejaVu Sans}
\setmonofont[Scale=0.82]{DejaVu Sans Mono}

\usepackage{amsmath,amssymb,amsthm,mathtools}
\usepackage{geometry}
\usepackage{graphicx}
\usepackage{booktabs,array,multirow,tabularx,longtable}
\usepackage{xcolor}
\usepackage[colorlinks=true,linkcolor=blue!65!black,citecolor=green!50!black,urlcolor=blue!70!black]{hyperref}
\usepackage{enumitem}
\usepackage{tikz}
\usetikzlibrary{positioning,arrows.meta,shapes.geometric,calc,
                decorations.pathreplacing,fit,backgrounds,matrix}
\usepackage{pgfplots}
\pgfplotsset{compat=1.18}
\usepackage{algorithm}
\usepackage{algpseudocode}
\usepackage{caption,subcaption}
\usepackage{fancyhdr}
\usepackage{titlesec}
\usepackage{tcolorbox}
\usepackage{listings}
\usepackage{float}
\usepackage{bm}
\usepackage{microtype}

\geometry{top=2.5cm,bottom=2.5cm,left=3cm,right=2.5cm}

\pagestyle{fancy}
\fancyhf{}
\rhead{\textcolor{gray!70}{\small RelGT++ -- Phân tích Kỹ thuật}}
\lhead{\textcolor{gray!70}{\small Từ mã nguồn đến toán học}}
\cfoot{\thepage}
\renewcommand{\headrulewidth}{0.3pt}

\titleformat{\section}{\Large\bfseries\color{blue!70!black}}{\thesection}{0.8em}{}[\color{blue!30!black}\titlerule]
\titleformat{\subsection}{\large\bfseries\color{blue!50!black}}{\thesubsection}{0.7em}{}
\titleformat{\subsubsection}{\normalsize\bfseries\color{black!75}}{\thesubsubsection}{0.6em}{}

\newtheorem{theorem}{Định lý}[section]
\newtheorem{proposition}[theorem]{Mệnh đề}
\newtheorem{definition}{Định nghĩa}[section]
\newtheorem{remark}{Nhận xét}[section]
\newtheorem{observation}{Quan sát}[section]

\tcbuselibrary{skins,breakable}

\newtcolorbox{codebox}[1][]{enhanced,colback=gray!5,colframe=gray!40,
  fonttitle=\bfseries\small\ttfamily,#1,breakable,left=4pt,right=4pt}

\newtcolorbox{mathbox}[1][]{enhanced,colback=green!4,colframe=green!55!black,
  fonttitle=\bfseries\small,#1,breakable}

\newtcolorbox{insightbox}[1][]{enhanced,colback=orange!4,colframe=orange!65!black,
  fonttitle=\bfseries\small,title=Nhận xét sâu,#1,breakable}

\newtcolorbox{issuebox}[1][]{enhanced,colback=red!4,colframe=red!50!black,
  fonttitle=\bfseries\small,title=Vấn đề kỹ thuật,#1,breakable}

\newtcolorbox{designbox}[1][]{enhanced,colback=blue!4,colframe=blue!55!black,
  fonttitle=\bfseries\small,title=Quyết định thiết kế,#1,breakable}

\lstset{
  basicstyle=\ttfamily\small,
  keywordstyle=\color{blue!70!black}\bfseries,
  commentstyle=\color{green!50!black}\itshape,
  stringstyle=\color{orange!70!black},
  backgroundcolor=\color{gray!6},
  frame=single,framerule=0.3pt,
  breaklines=true,
  numbers=left,numberstyle=\tiny\color{gray},
  numbersep=5pt,
  xleftmargin=15pt,
  language=Python
}

% ---------------------------------------------------------------
\begin{document}

% ==================== TRANG BÌA ====================
\begin{titlepage}
\centering\vspace*{1cm}

{\small\scshape\color{gray} Phân tích Kỹ thuật -- Từ Mã nguồn đến Toán học}\\[1em]
\rule{\textwidth}{2.5pt}\\[0.7em]
{\Huge\bfseries\color{blue!75!black} Phân tích Chuyên sâu\\[0.4em]
\Large Kiến trúc \texttt{RelGT++}
}\\[0.5em]
\rule{\textwidth}{1pt}\\[2em]

\begin{tcolorbox}[colback=blue!4,colframe=blue!50,width=0.92\textwidth,
  title={\bfseries Phạm vi phân tích},fonttitle=\bfseries]
Tài liệu này phân tích \textbf{toàn bộ mã nguồn RelGT++} theo từng module,
truy xuất công thức toán học trực tiếp từ code, xác định các quyết định
thiết kế, và chỉ rõ điểm mạnh, điểm yếu, cùng các vấn đề kỹ thuật tiềm ẩn.

\medskip
\textbf{Các file được phân tích:}
\texttt{model.py} $\cdot$ \texttt{codebook.py} $\cdot$
\texttt{local\_module.py} $\cdot$ \texttt{encoders.py} $\cdot$
\texttt{utils.py} $\cdot$ \texttt{main\_node\_ddp.py}
\end{tcolorbox}

\vspace{2em}
\begin{tabular}{lp{9cm}}
\toprule
\textbf{Kiến trúc gốc} & RelGT (Dwivedi et al., arXiv:2505.10960) \\
\textbf{Phiên bản phân tích} & RelGT++ (triển khai cải tiến) \\
\textbf{Tổng số module mới} & 6 module (so với RelGT gốc) \\
\textbf{Tổng dòng code} & $\approx$1800 dòng (không tính utils) \\
\textbf{Ngày phân tích} & \today \\
\bottomrule
\end{tabular}
\vfill
\end{titlepage}

\tableofcontents\newpage

% ======================================================================
\section{Bản đồ Kiến trúc Tổng thể}
% ======================================================================

\subsection{Luồng dữ liệu đầy đủ}

Trước khi đi vào từng module, ta cần nắm rõ luồng tensor qua toàn bộ hệ thống.

\begin{figure}[H]
\centering
\begin{tikzpicture}[
  node distance=0.45cm and 0.9cm,
  box/.style={draw,rounded corners=4pt,minimum width=2.6cm,minimum height=0.75cm,
              align=center,font=\footnotesize},
  nb/.style={box,fill=blue!12,draw=blue!55!black},
  ob/.style={box,fill=gray!12,draw=gray!45},
  sb/.style={box,fill=green!10,draw=green!55!black},
  ab/.style={box,fill=orange!10,draw=orange!55!black},
  arr/.style={-Stealth,thick},
  dim/.style={font=\tiny\color{red!65!black},align=center}
]

%% Column 1: Input
\node[ob] (inp_nt) {neighbor\_types};
\node[ob,below=0.35cm of inp_nt] (inp_nh) {neighbor\_hops};
\node[ob,below=0.35cm of inp_nh] (inp_ni) {neighbor\_times};
\node[ob,below=0.35cm of inp_ni] (inp_tf) {grouped\_tfs};
\node[ob,below=0.35cm of inp_tf] (inp_ei) {edge\_index};

%% Column 2: Encoders
\node[nb,right=1.1cm of inp_nt] (enc_t)  {TypeEncoder};
\node[nb,right=1.1cm of inp_nh] (enc_h)  {HopEncoder};
\node[nb,right=1.1cm of inp_ni] (enc_ti) {TimeEncoder};
\node[nb,right=1.1cm of inp_tf] (enc_f)  {TfsEncoder\\(ResNet/TF)};
\node[nb,right=1.1cm of inp_ei] (enc_pe) {GNNPEEncoder\\(GIN x4)};

%% Dims after encoding
\node[dim,right=0.15cm of enc_t]  {[B,K,d]};
\node[dim,right=0.15cm of enc_h]  {[B,K,d]};
\node[dim,right=0.15cm of enc_ti] {[B,K,d]};
\node[dim,right=0.15cm of enc_f]  {[B,K,d]};
\node[dim,right=0.15cm of enc_pe] {[B,K,d]};

%% Column 3: Fusion
\node[ab,right=2.5cm of enc_h,yshift=-0.5cm]  (fusion) {CrossModal\\GatedFusion};
\node[dim,right=0.15cm of fusion] {[B,K,d]};

%% Column 4: RelGT Layer
\node[nb,right=1.3cm of fusion,yshift=1.2cm]  (local)  {LocalModule\\(Transformer)};
\node[nb,right=1.3cm of fusion,yshift=-1.2cm] (global) {GlobalModule\\(VQ Attn)};
\node[dim,right=0.1cm of local]  {[B,d]};
\node[dim,right=0.1cm of global] {[B,d]};

%% Column 5: Bridge + FF
\node[sb,right=1.2cm of local,yshift=-1.2cm]  (bridge) {CrossAttn\\Bridge};
\node[nb,right=1.2cm of bridge] (ff) {FF + BN};
\node[dim,right=0.1cm of ff] {[B,d]};

%% Column 6: Head
\node[ab,right=1.1cm of ff] (head) {MLP Head};
\node[ob,right=1cm of head] (out)  {$\hat{y}$};

%% Arrows: inputs to encoders
\foreach \s/\t in {inp_nt/enc_t,inp_nh/enc_h,inp_ni/enc_ti,inp_tf/enc_f,inp_ei/enc_pe}
  \draw[arr,gray!60] (\s) -- (\t);

%% Encoders to fusion
\foreach \s in {enc_t,enc_h,enc_ti,enc_f,enc_pe}
  \draw[arr,blue!50] (\s.east) -- ++(0.35,0) |- (fusion.west);

%% Fusion to local and global
\draw[arr] (fusion.east) -- ++(0.4,0) |- (local.west);
\draw[arr] (fusion.east) -- ++(0.4,0) |- (global.west)
  node[midway,right,font=\tiny]{$x=x\_set[:,0,:]$};

%% Local + Global to bridge
\draw[arr] (local.east)  -- ++(0.3,0) |- (bridge.west);
\draw[arr] (global.east) -- ++(0.3,0) |- (bridge.west);

%% Bridge to FF
\draw[arr] (bridge) -- (ff);

%% FF to head to output
\draw[arr] (ff) -- (head);
\draw[arr] (head) -- (out);

%% Aux loss annotation
\node[font=\tiny\color{red!60},below=0.15cm of bridge] {aux losses};

\end{tikzpicture}
\caption{Luồng tensor đầy đủ trong RelGT++. Số liệu chiều tensor được ghi màu đỏ.}
\label{fig:full_flow}
\end{figure}

\subsection{Bảng tham số theo module}

\begin{table}[H]
\centering\small
\caption{Kích thước tham số ước tính ($d=512$, $B_\text{centroids}=4096$, $K=300$, $H=4$)}
\begin{tabular}{llrr}
\toprule
\textbf{Module} & \textbf{File} & \textbf{Tham số (approx)} & \textbf{\% tổng} \\
\midrule
NeighborTfsEncoder (ResNet/TF) & encoders.py & $\sim$8.4M & 54\% \\
CrossModalGatedFusion & model.py & $\sim$3.9M & 25\% \\
LocalModule (Transformer) & local\_module.py & $\sim$1.6M & 10\% \\
GlobalModule (VQ Attn) & model.py & $\sim$0.5M & 3\% \\
GumbelSoftCodebook & codebook.py & $\sim$0.4M & 2\% \\
CrossAttentionBridge & model.py & $\sim$0.3M & 2\% \\
GNNPEEncoder (GIN) & encoders.py & $\sim$0.2M & 1\% \\
TypeEncoder + HopEncoder & encoders.py & $\sim$0.05M & $<$1\% \\
TimeEncoder & encoders.py & $\sim$0.05M & $<$1\% \\
MLP Head & model.py & $\sim$0.26M & 2\% \\
\midrule
\textbf{Tổng} & & $\sim$\textbf{15.6M} & 100\% \\
\bottomrule
\end{tabular}
\end{table}

% ======================================================================
\section{Module 1: Hệ thống Encoder (encoders.py)}
% ======================================================================

\subsection{NeighborNodeTypeEncoder}

\subsubsection{Mã nguồn và công thức}

\begin{codebox}[title={encoders.py -- NeighborNodeTypeEncoder.forward()}]
\begin{lstlisting}
def forward(self, type_indices):
    return self.embedding(type_indices)
# nn.Embedding(num_types + 1, embedding_dim)
\end{lstlisting}
\end{codebox}

\begin{mathbox}[title={Công thức toán học}]
\begin{equation}
  h_\text{type}(v_j) = W_\text{type}[\varphi(v_j)] \in \mathbb{R}^d
  \label{eq:type_enc}
\end{equation}
với $W_\text{type} \in \mathbb{R}^{(|\mathcal{T}|+1) \times d}$ là ma trận nhúng.
Chiều $+1$ để xử lý kiểu ``unknown''. Kích thước từ điển thực:
$\texttt{num\_types} = \max(\text{node\_type\_map.values()}) + 1$.
\end{mathbox}

\begin{insightbox}
\textbf{Khác biệt so với RelGT gốc:} RelGT gốc sử dụng one-hot rồi nhân với ma trận
$W_\text{type}$ tường minh (công thức~\ref{eq:type_enc} trong paper). RelGT++ dùng
\texttt{nn.Embedding} -- về mặt toán học tương đương, nhưng hiệu quả hơn vì tránh
phép nhân ma trận one-hot thưa thớt.
\end{insightbox}

\subsection{NeighborHopEncoder}

\begin{codebox}[title={encoders.py -- NeighborHopEncoder.forward()}]
\begin{lstlisting}
def forward(self, hop_distances):
    shifted = hop_distances + 1   # shift: hop 1 -> idx 2, hop 0 -> idx 1
    return self.embedding(shifted)
# nn.Embedding(max_neighbor_hop + 2, embedding_dim)
\end{lstlisting}
\end{codebox}

\begin{mathbox}[title={Công thức toán học}]
\begin{equation}
  h_\text{hop}(v_i, v_j) = W_\text{hop}[p(v_i,v_j) + 1] \in \mathbb{R}^d
\end{equation}
với $p(v_i,v_j) \in \{0,1,2,3\}$, $W_\text{hop} \in \mathbb{R}^{(h_\text{max}+2)\times d}$.
Dịch chuyển $+1$ để chỉ số 0 giữ nguyên cho padding/unknown.
\end{mathbox}

\begin{designbox}
\textbf{max\_neighbor\_hop = 3} trong RelGTTokens (utils.py, dòng 234):
\begin{center}
\texttt{self.max\_neighbor\_hop = 2 + 1}
\end{center}
Nghĩa là: hop=1 (lân cận 1 bước), hop=2 (lân cận 2 bước), hop=3 (fallback ngẫu nhiên).
Từ điển embedding có kích thước $3+2=5$.
\end{designbox}

\subsection{NeighborTimeEncoder}

\subsubsection{Thiết kế hai giai đoạn}

\begin{codebox}[title={encoders.py -- NeighborTimeEncoder.forward() (rút gọn)}]
\begin{lstlisting}
def forward(self, rel_time):   # rel_time: [B, K] seconds
    B, K = rel_time.shape
    flattened_time = rel_time.view(-1)              # [B*K]
    pos_encoded = self.pos_encoder(flattened_time)  # [B*K, d]
    linear_out  = self.linear(pos_encoded)          # [B*K, d]
    linear_out  = linear_out.view(B, K, -1)         # [B, K, d]
    mask        = (rel_time < 0).unsqueeze(-1).float()
    out = (1 - mask) * linear_out + mask * self.mask_vector
    return out
\end{lstlisting}
\end{codebox}

\begin{mathbox}[title={Công thức toán học đầy đủ}]
\textbf{Bước 1 -- Positional Encoding (Vaswani et al.):}
\begin{equation}
  \text{PE}(\Delta\tau)_{2k}   = \sin\!\left(\frac{\Delta\tau}{10000^{2k/d}}\right),
  \quad
  \text{PE}(\Delta\tau)_{2k+1} = \cos\!\left(\frac{\Delta\tau}{10000^{2k/d}}\right)
  \label{eq:pe_sin}
\end{equation}

\textbf{Bước 2 -- Chiếu tuyến tính:}
\begin{equation}
  h'_\text{time} = W_\text{time} \cdot \text{PE}(\Delta\tau) + b_\text{time}
  \in \mathbb{R}^d
\end{equation}

\textbf{Bước 3 -- Xử lý token không có thời gian ($\Delta\tau < 0$):}
\begin{equation}
  \boxed{h_\text{time}(v_j) = (1 - m_j) \cdot h'_\text{time} + m_j \cdot \bm{\mu}_\text{mask}}
  \label{eq:time_enc}
\end{equation}
với $m_j = \mathbf{1}[\Delta\tau_j < 0]$ và $\bm{\mu}_\text{mask} \in \mathbb{R}^d$
là tham số có thể học khởi tạo bằng $\mathbf{0}$.
\end{mathbox}

\begin{insightbox}
\textbf{Tại sao dùng giây ($\Delta\tau$ tính bằng giây)?} Trong \texttt{utils.py}:
\begin{center}
\texttt{relative\_time\_days = (seed\_time - nbr\_time) / (60 * 60 * 24)}
\end{center}
Dữ liệu lưu bằng \emph{ngày}, nhưng \texttt{NeighborTimeEncoder} trong
\texttt{model.py} gọi \texttt{neighbor\_times.float()} trực tiếp -- không có
chuyển đổi đơn vị bổ sung. Giá trị $\Delta\tau$ đưa vào PE là số ngày.
PE có chu kỳ $2\pi \cdot 10000^{2k/d}$ -- với $d=512$ và $k=0$, chu kỳ là
$2\pi \approx 6.28$ ngày. Điều này nghĩa là PE có độ phân giải tốt cho
khoảng thời gian $\sim$1 tuần trở xuống.
\end{insightbox}

\subsection{NeighborTfsEncoder (Encoder chính -- 54\% tham số)}

\subsubsection{Kiến trúc ResNet đa phương thức}

\begin{mathbox}[title={Pipeline mã hóa đặc trưng bảng}]
Với mỗi loại node $t \in \mathcal{T}$ và batch $\mathbf{X}^{(t)}$
gồm $N_t$ hàng từ bảng $t$:

\textbf{Bước 1 -- Mã hóa theo phương thức:}
\begin{align}
  h_\text{cat}^{(t)}  &= \text{EmbeddingEncoder}(\mathbf{X}^{(t)}_\text{cat})
    & &(\text{phân loại})\\
  h_\text{num}^{(t)}  &= \text{LinearEncoder}(\mathbf{X}^{(t)}_\text{num})
    & &(\text{số})\\
  h_\text{text}^{(t)} &= \text{LinearEmbeddingEncoder}(\mathbf{X}^{(t)}_\text{text})
    & &(\text{văn bản -- GloVe})\\
  h_\text{ts}^{(t)}   &= \text{TimestampEncoder}(\mathbf{X}^{(t)}_\text{ts})
    & &(\text{nhãn thời gian})
\end{align}

\textbf{Bước 2 -- ResNet tổng hợp:}
\begin{equation}
  h_\text{feat}^{(t)} = \text{ResNet}\!\left(
    \bigoplus_{m} h_m^{(t)};\;
    \theta^{(t)}_\text{resnet}
  \right) \in \mathbb{R}^{N_t \times d}
  \label{eq:tfs_enc}
\end{equation}
với $\bigoplus$ là phép gộp qua các phương thức (mean/concat trong TF ResNet).
\end{mathbox}

\subsubsection{Chiến lược Grouping và Scatter}

Đây là điểm kỹ thuật quan trọng nhất trong NeighborTfsEncoder:

\begin{codebox}[title={encoders.py -- forward() -- nhóm theo kiểu node}]
\begin{lstlisting}
# Pre-allocate flat buffer [N_total, d]
encoded_flat_tensor = torch.zeros((N, self.channels), device=device)

for t_int, big_tf in grouped_tfs.items():        # one pass per node type
    out_t = encoder(big_tf)                       # [N_t, d]
    idx_tensor = torch.tensor(grouped_indices[t_int], ...)
    encoded_flat_tensor[idx_tensor] = out_t       # scatter back

# Scatter [N_total, d] -> [B, K, d]
output[indices_i, indices_j] = encoded_flat_tensor
\end{lstlisting}
\end{codebox}

\begin{mathbox}[title={Sơ đồ Grouping-Encode-Scatter}]
\begin{equation}
  \underbrace{[B, K]}_{\text{neighbor\_types}}
  \xrightarrow{\text{group by type}}
  \underbrace{\{(t, \text{big\_tf}_t)\}_{t \in \mathcal{T}}}_{\text{nhóm theo bảng}}
  \xrightarrow{\text{ResNet per type}}
  \underbrace{[N_\text{total}, d]}_{\text{flat encodings}}
  \xrightarrow{\text{scatter}}
  \underbrace{[B, K, d]}_{\text{output}}
\end{equation}
\end{mathbox}

\begin{insightbox}
\textbf{Tại sao grouping quan trọng?} Nếu mã hóa từng token riêng lẻ,
ta cần $B \times K$ lần forward qua ResNet. Với $B=512$, $K=300$: \textbf{153,600 lần}.
Với grouping, chỉ cần $|\mathcal{T}|$ lần ($\sim$3--8 lần tùy dataset).
Speedup lý thuyết: $\sim$20,000--50,000$\times$.
\end{insightbox}

\begin{issuebox}
\textbf{NaN/Inf handling (dòng 233--237):}
\begin{lstlisting}[numbers=none]
big_tf.feat_dict[stype] = torch.nan_to_num(tensor, nan=0.0, posinf=1e6, neginf=-1e6)
\end{lstlisting}
Đây là biện pháp phòng thủ nhưng có thể che giấu lỗi dữ liệu. Giá trị $10^6$
được clip vào đặc trưng số học có thể tạo ra gradient bùng nổ nếu không được
normalize trước. \textbf{Khuyến nghị:} Thêm kiểm tra cảnh báo (warning) khi
tỉ lệ NaN $>1\%$.
\end{issuebox}

\subsection{GNNPEEncoder}

\subsubsection{Kiến trúc GIN với Residual}

\begin{codebox}[title={encoders.py -- GNNPEEncoder.forward() (rút gọn)}]
\begin{lstlisting}
x_input = torch.randn(total_nodes, 1, device=device)  # stochastic init
x = self.input_proj(x_input)                           # [N, d//4]

outputs = []
for i, conv in enumerate(self.conv):           # 4 GIN layers
    x_res = x
    x_new = conv(x, edge_index)                # GINConv
    x_new = self.bns[i](x_new)
    x_new = F.relu(x_new)
    x = x_new + x_res                          # residual connection
    outputs.append(x)

x = self.final_transform(outputs[-1])          # [N, d]
out = x.view(B, K, -1)                         # [B, K, d]
\end{lstlisting}
\end{codebox}

\begin{mathbox}[title={Công thức GIN với Residual}]
\textbf{Khởi tạo ngẫu nhiên:} $Z^{(0)} \sim \mathcal{N}(0, I) \in \mathbb{R}^{K \times 1}$

\textbf{Input projection:} $X^{(0)} = W_\text{proj} Z^{(0)} \in \mathbb{R}^{K \times d/4}$

\textbf{GIN layer $l$:}
\begin{equation}
  X^{(l)}_\text{new} = \text{BN}\!\left(
    \text{MLP}_l\!\left(
      (1 + \epsilon_l) X^{(l-1)} + \sum_{j \in \mathcal{N}(v)} X^{(l-1)}_j
    \right)
  \right)
\end{equation}
\begin{equation}
  X^{(l)} = \text{ReLU}(X^{(l)}_\text{new}) + X^{(l-1)}
  \quad \text{(residual)}
  \label{eq:gin_res}
\end{equation}

\textbf{Chiếu cuối:}
\begin{equation}
  h_\text{pe}(v_j) = W_\text{final} X^{(L)}_j \in \mathbb{R}^d
\end{equation}
\end{mathbox}

\begin{observation}
\textbf{Chiều ẩn = $d/4$:}
\begin{center}
\texttt{self.layer\_embedding\_dim = embedding\_dim // 4}
\end{center}
Với $d=512$, GIN hoạt động ở $d_\text{GIN}=128$, rồi chiếu lên $d=512$ ở bước cuối.
Thiết kế này tiết kiệm $\sim$16$\times$ tham số trong các lớp GIN
($4 \times 128^2 = 65\text{K}$ thay vì $4 \times 512^2 = 1\text{M}$).
\end{observation}

\begin{issuebox}
\textbf{Permutation equivariance bị vi phạm:} Dòng 346:
\begin{center}
\texttt{x\_input = torch.randn(total\_nodes, 1, device=device)}
\end{center}
Mỗi lần forward tạo ra $Z_\text{random}$ khác nhau. Điều này có nghĩa:
với cùng một đồ thị con, hai lần chạy cho ra $h_\text{pe}$ khác nhau.
Đây là \emph{intentional} (theo paper, giúp tăng expressiveness) nhưng
gây ra \textbf{variance cao trong inference} -- mô hình phải học cách
trung bình hóa qua nhiều lần random initialization.

\textbf{Vấn đề thực tế:} Khi $K=300$ node và $L=4$ lớp GIN, gradient
của $Z_\text{random}$ không tồn tại (detached bởi \texttt{randn}).
Chỉ tham số $W_\text{proj}$ và $W_\text{final}$ nhận gradient -- nghĩa là
GIN học cách ``đọc cấu trúc graph'' từ nhiễu ngẫu nhiên, không phải từ
đặc trưng xác định.
\end{issuebox}

% ======================================================================
\section{Module 2: CrossModalGatedFusion (model.py)}
% ======================================================================

\subsection{Phân tích Công thức}

\begin{codebox}[title={model.py -- CrossModalGatedFusion.forward()}]
\begin{lstlisting}
def forward(self, embeddings: list):   # list of N tensors [B, K, C]
    total = sum(embeddings)            # [B, K, C]  -- sum of all N modalities
    fused = torch.zeros_like(embeddings[0])
    for i in range(N):
        ctx = (total - embeddings[i]) / max(N - 1, 1)  # mean of others
        gate = torch.sigmoid(
            self.self_projs[i](embeddings[i]) + self.ctx_projs[i](ctx)
        )
        val = self.val_projs[i](embeddings[i])
        fused = fused + gate * val
    return self.layer_norm(self.out_proj(fused))
\end{lstlisting}
\end{codebox}

\begin{mathbox}[title={Công thức CrossModalGatedFusion đầy đủ}]
Với $N=5$ phương thức, mỗi $e_i \in \mathbb{R}^{B \times K \times d}$:

\textbf{Bước 1 -- Context cho phương thức $i$:}
\begin{equation}
  \text{ctx}_i = \frac{1}{N-1}\sum_{j \ne i} e_j
  = \frac{\sum_{j=1}^N e_j - e_i}{N-1}
  \label{eq:cmgf_ctx}
\end{equation}

\textbf{Bước 2 -- Gate:}
\begin{equation}
  g_i = \sigma\!\left(W^i_\text{self}\, e_i + W^i_\text{ctx}\, \text{ctx}_i\right)
  \in [0,1]^{B \times K \times d}
  \label{eq:cmgf_gate}
\end{equation}

\textbf{Bước 3 -- Value:}
\begin{equation}
  \text{val}_i = W^i_\text{val}\, e_i \in \mathbb{R}^{B \times K \times d}
\end{equation}

\textbf{Bước 4 -- Fusion:}
\begin{equation}
  f = \sum_{i=1}^N g_i \odot \text{val}_i \in \mathbb{R}^{B \times K \times d}
\end{equation}

\textbf{Bước 5 -- Output:}
\begin{equation}
  \boxed{h_\text{fused} = \text{LN}\!\left(W_\text{out}\, f\right) \in \mathbb{R}^{B \times K \times d}}
  \label{eq:cmgf_out}
\end{equation}
\end{mathbox}

\subsection{So sánh với RelGT gốc}

\begin{table}[H]
\centering\small
\caption{So sánh cơ chế kết hợp phương thức}
\begin{tabular}{lll}
\toprule
\textbf{Tính chất} & \textbf{RelGT gốc} & \textbf{RelGT++ (CMGF)} \\
\midrule
Công thức & $h = W_O \cdot [e_1\|e_2\|e_3\|e_4\|e_5]$ & $h = \text{LN}(W_O \sum_i g_i \odot W_i^v e_i)$ \\
Tham số & $W_O \in \mathbb{R}^{5d \times d}$ & $3Nd^2 + d^2$ ($\approx 15d^2 + d^2$) \\
Phương thức tắt & Không thể & Có thể: $g_i \to 0$ \\
Per-token adaptive & Không & Có \\
Gradient qua phương thức & Đồng đều & Được điều chỉnh bởi gate \\
\bottomrule
\end{tabular}
\end{table}

\begin{insightbox}
\textbf{Lợi thế chính:} Gate $g_i$ phụ thuộc vào cả nội dung của phương thức $i$
\emph{và} bối cảnh từ tất cả phương thức khác. Điều này cho phép mô hình học rằng,
ví dụ, \emph{khi đặc trưng bảng của node rất giàu thông tin} (giá trị $e_4$ lớn),
hãy giảm tầm quan trọng của hop encoder (giảm $g_3$).

\textbf{Số tham số:} $N=5$ phương thức, mỗi phương thức có 3 ma trận $d \times d$
($W_\text{self}$, $W_\text{ctx}$, $W_\text{val}$) cộng thêm $W_\text{out}$:
$(3 \times 5 + 1) \times 512^2 = 4{,}194{,}304 \approx 4$M tham số.
\end{insightbox}

\begin{issuebox}
\textbf{Vấn đề gradient cho phương thức yếu:} Nếu phương thức $i$ không mang
thông tin, gate $g_i \to 0$ -- nhưng gradient cũng bị chặn tại đây:
$\partial \mathcal{L} / \partial e_i \approx g_i \cdot (\ldots) \approx 0$.
Điều này có thể dẫn đến tình trạng một phương thức ``chết'' hoàn toàn
(dead modality) và không thể phục hồi sau đó.
\end{issuebox}

% ======================================================================
\section{Module 3: LocalModule và EncoderLayer (local\_module.py)}
% ======================================================================

\subsection{Kiến trúc Tổng thể LocalModule}

\begin{mathbox}[title={LocalModule forward pass}]
\begin{equation}
  X^{(0)} = W_\text{embed}\, h_\text{fused} \in \mathbb{R}^{B \times K \times d}
  \quad (d_\text{in} \to d_\text{hidden})
\end{equation}
\begin{equation}
  X^{(l)} = \text{EncoderLayer}_l(X^{(l-1)}), \quad l = 1, \ldots, L
\end{equation}
\begin{equation}
  X_\text{final} = \text{LN}(X^{(L)}) \in \mathbb{R}^{B \times K \times d}
\end{equation}
\begin{equation}
  h_\text{local}(v_i) = \text{MultiViewReadout}(X_\text{final}) \in \mathbb{R}^{B \times d}
\end{equation}
\end{mathbox}

\subsection{StochasticDepth}

\begin{codebox}[title={local\_module.py -- StochasticDepth.forward()}]
\begin{lstlisting}
def forward(self, x):
    if not self.training or self.drop_prob == 0.0:
        return x
    keep = torch.rand(1, device=x.device) > self.drop_prob
    return x * keep.float() / (1.0 - self.drop_prob)
\end{lstlisting}
\end{codebox}

\begin{mathbox}[title={Công thức Stochastic Depth (DropPath)}]
\begin{equation}
  \text{SD}(x;\, p) = \frac{x \cdot \mathbf{1}[\text{Bernoulli}(1-p)]}{1-p}
  \label{eq:sd}
\end{equation}

Tỉ lệ drop tăng tuyến tính theo chiều sâu lớp:
\begin{equation}
  p_l = 0.05 \cdot \frac{l+1}{\max(L, 1)}, \quad l = 0, 1, \ldots, L-1
  \label{eq:sd_rate}
\end{equation}
Với $L=1$: $p_0 = 0.05$.
\end{mathbox}

\begin{observation}
\textbf{DropPath vs Dropout:} Dropout loại bỏ từng \emph{phần tử} độc lập.
DropPath (Stochastic Depth) loại bỏ toàn bộ \emph{layer output} cho một
sample trong batch. Điều này hiệu quả hơn cho Transformer vì nó buộc
residual connection phải học biểu diễn đủ tốt mà không cần sự trợ giúp
từ sub-layer.

\textbf{Lưu ý quan trọng:} Trong code, \texttt{torch.rand(1)} tạo
\emph{một giá trị scalar} -- tức toàn bộ batch cùng drop hay không.
Đây là \textbf{layer-level} stochastic depth, không phải sample-level.
\end{observation}

\subsection{GatedFFN (SwiGLU)}

\begin{codebox}[title={local\_module.py -- GatedFFN.forward()}]
\begin{lstlisting}
def forward(self, x):         # x: [B, K, D]
    x = permute_bn(x)         # BN over D dimension
    gate = F.silu(self.w_gate(x))    # [B, K, ffn_size]
    up   = self.w_up(x)              # [B, K, ffn_size]
    x    = self.w_down(self.dropout(gate * up))  # [B, K, D]
    x    = self.dropout(x)
    x    = permute_bn_out(x)  # BN over D
    return x
\end{lstlisting}
\end{codebox}

\begin{mathbox}[title={Công thức SwiGLU FFN}]
\begin{equation}
  \text{GatedFFN}(x) = \text{BN}_\text{out}\!\left(
    W_\text{down} \bigl(
      \underbrace{\text{SiLU}(W_\text{gate}\, x)}_{\text{cổng}} \odot
      \underbrace{W_\text{up}\, x}_{\text{giá trị}}
    \bigr)
  \right)
  \label{eq:swiglu}
\end{equation}
với $\text{SiLU}(x) = x \cdot \sigma(x)$ (Sigmoid Linear Unit).

\textbf{Kích thước:}
\begin{itemize}[nosep]
  \item $W_\text{gate}, W_\text{up} \in \mathbb{R}^{d \times d_\text{ffn}}$
  \item $W_\text{down} \in \mathbb{R}^{d_\text{ffn} \times d}$
  \item $d_\text{ffn} = \texttt{ffn\_size} \times 2 = 2d_\text{hidden}$
        (vì EncoderLayer truyền \texttt{ffn\_size * 2})
\end{itemize}
Thực tế: LocalModule đặt \texttt{ffn\_dim = 2 * hidden\_dim}, rồi
EncoderLayer truyền vào GatedFFN với \texttt{ffn\_size * 2}.
Vậy expansion factor thực là $\mathbf{4\times}$.
\end{mathbox}

\begin{insightbox}
\textbf{BatchNorm trên 3D tensor:} Code dùng \texttt{x.permute(0, 2, 1)}
trước khi áp dụng \texttt{nn.BatchNorm1d}. Lý do: BN1d mong đợi
\texttt{[N, C, L]} (channel ở dim 1), nên cần permute từ \texttt{[B, K, D]}
sang \texttt{[B, D, K]}, rồi permute ngược lại.

\textbf{Hàm ý thực tế:} BN tính thống kê trên chiều $K$ (tức là trên tất cả
$K=300$ token trong cùng một sample). Điều này khác BN trên batch ($B$) --
batch BN ở đây thực ra là \emph{feature normalization theo không gian token}.
\end{insightbox}

\subsection{EncoderLayer: Attention + DropPath}

\begin{mathbox}[title={EncoderLayer forward pass đầy đủ}]
\textbf{Self-Attention với DropPath:}
\begin{align}
  Q, K, V &= W_q X_\text{norm},\; W_k X_\text{norm},\; W_v X_\text{norm} \\
  A &= \text{FlashAttn}(Q, K, V;\; \text{dropout}=p_\text{attn}) \\
  A &= W_\text{out} A \\
  X &\leftarrow X + \text{SD}(A;\; p_l) \label{eq:enc_attn}
\end{align}

\textbf{FFN với DropPath:}
\begin{align}
  F &= \text{GatedFFN}(\text{LN}(X)) \\
  X &\leftarrow X + \text{SD}(F;\; p_l) \label{eq:enc_ffn}
\end{align}

\textbf{Flash Attention} (PyTorch 2.0+):
\begin{equation}
  \text{FlashAttn}(Q, K, V) = \text{softmax}\!\left(\frac{QK^\top}{\sqrt{d_k}}\right) V
  \quad \text{(IO-aware implementation)}
\end{equation}
\end{mathbox}

\begin{issuebox}
\textbf{DropPath áp dụng hai lần cùng giá trị:} Trong forward của EncoderLayer,
cùng một \texttt{self.drop\_path} được gọi cho cả attention và FFN output.
Vì \texttt{torch.rand(1)} được gọi riêng trong mỗi lần forward của StochasticDepth,
hai lần drop là \textbf{độc lập} -- đây là thiết kế đúng.

Tuy nhiên, nếu \texttt{drop\_path\_rate = 0.0}, code trả về \texttt{nn.Identity()}
không phải \texttt{StochasticDepth(0.0)} -- điều này tránh overhead của một
module gọi không cần thiết. Thiết kế tốt.
\end{issuebox}

\subsection{MultiViewReadout}

\begin{codebox}[title={local\_module.py -- MultiViewReadout.forward()}]
\begin{lstlisting}
node_tensor     = output[:, 0, :]        # [B, D] seed node
neighbor_tensor = output[:, 1:, :]       # [B, K-1, D]

# View 1: attention-weighted
target      = node_tensor.unsqueeze(1).expand_as(neighbor_tensor)
attn_scores = self.attn_layer(torch.cat([target, neighbor_tensor], dim=-1))
attn_weights = F.softmax(attn_scores, dim=1)    # softmax over K-1 neighbors
v_attn = (neighbor_tensor * attn_weights).sum(dim=1)   # [B, D]

# View 2, 3
v_mean = neighbor_tensor.mean(dim=1)             # [B, D]
v_max  = neighbor_tensor.max(dim=1)[0]           # [B, D]

combined = torch.cat([v_attn, v_mean, v_max], dim=-1)  # [B, 3D]
fused    = self.combine(combined)                        # [B, D]
return self.layer_norm(node_tensor + fused)              # residual
\end{lstlisting}
\end{codebox}

\begin{mathbox}[title={Công thức Multi-View Readout}]
\textbf{View 1 -- Attention seed-to-neighbor:}
\begin{equation}
  \alpha_j = \frac{\exp(w^\top [h_0 \| h_j])}{\sum_{k=1}^{K-1} \exp(w^\top [h_0 \| h_k])},
  \quad v_\text{attn} = \sum_{j=1}^{K-1} \alpha_j h_j
  \label{eq:mvr_attn}
\end{equation}

\textbf{View 2, 3 -- Aggregation:}
\begin{equation}
  v_\text{mean} = \frac{1}{K-1}\sum_{j=1}^{K-1} h_j, \qquad
  v_\text{max}  = \max_{j=1}^{K-1} h_j
  \label{eq:mvr_pool}
\end{equation}

\textbf{Combine với residual:}
\begin{equation}
  \boxed{h_\text{local}(v_0) = \text{LN}\!\left(
    h_0 + \text{MLP}\!\left([v_\text{attn} \| v_\text{mean} \| v_\text{max}]\right)
  \right)}
  \label{eq:mvr_out}
\end{equation}
\end{mathbox}

\begin{insightbox}
\textbf{Ba góc nhìn bổ sung cho nhau:}
\begin{itemize}[nosep]
  \item $v_\text{attn}$: ``Lân cận nào quan trọng với seed node?'' -- context-aware
  \item $v_\text{mean}$: ``Đặc trưng trung bình của toàn bộ vùng lân cận'' -- global stats
  \item $v_\text{max}$: ``Tính chất nổi bật nhất trong vùng lân cận'' -- salient feature
\end{itemize}
MLP$(3d \to 2d \to d)$ học cách kết hợp ba góc nhìn này một cách tối ưu.
\end{insightbox}

% ======================================================================
\section{Module 4: GumbelSoftCodebook (codebook.py)}
% ======================================================================

\subsection{Vấn đề Codebook Collapse}

Trong VQ-EMA gốc, cập nhật codebook dùng EMA:
\begin{equation}
  e_b \leftarrow \frac{\sum_{v:\,c_v=b} q_v}{\sum_{v:\,c_v=b} 1}
\end{equation}
Nếu centroid $b$ không được sử dụng trong nhiều batch, $N_b \to 0$ và
$e_b$ không được cập nhật -- dẫn đến collapse.

\subsection{Cơ chế Gumbel-Softmax}

\begin{codebox}[title={codebook.py -- GumbelSoftCodebook.update()}]
\begin{lstlisting}
distances = (x_norm.pow(2).sum(1,keepdim=True)
             - 2 * x_norm @ self.embedding_k.t()
             + self.embedding_k.pow(2).sum(1).unsqueeze(0))
logits = -distances

if self.training:
    gumbels = -torch.log(-torch.log(torch.rand_like(logits) + 1e-20) + 1e-20)
    noisy_logits = (logits + gumbels) / self.temperature.clamp(min=0.1)
    soft_probs   = F.softmax(noisy_logits, dim=-1)
    hard_indices = soft_probs.argmax(dim=-1)       # straight-through
    
    # EMA usage tracking
    counts = torch.bincount(hard_indices, minlength=self._num_embeddings).float()
    counts = counts / (counts.sum() + 1e-8)
    self._ema_usage.mul_(0.99).add_(counts * 0.01)
    
    # Temperature annealing
    new_temp = max(self.temperature.item() * (1 - self.anneal_rate), self.min_temp)
    self.temperature.fill_(new_temp)
\end{lstlisting}
\end{codebox}

\begin{mathbox}[title={Công thức GumbelSoftCodebook đầy đủ}]
\textbf{Distance-based logits:}
\begin{equation}
  l_{v,b} = -\|q_v - e_b\|_2^2
    = -\|q_v\|^2 + 2\langle q_v, e_b\rangle - \|e_b\|^2
  \label{eq:gsb_logit}
\end{equation}

\textbf{Gumbel noise:}
\begin{equation}
  g_{v,b} \sim \text{Gumbel}(0,1) = -\log(-\log U_{v,b}),\quad U_{v,b} \sim \text{Uniform}(0,1)
\end{equation}

\textbf{Soft assignment:}
\begin{equation}
  s_{v,b} = \frac{\exp\bigl((l_{v,b} + g_{v,b})/\tau\bigr)}
    {\sum_{b'}\exp\bigl((l_{v,b'} + g_{v,b'})/\tau\bigr)}
  \label{eq:gsb_soft}
\end{equation}

\textbf{Hard assignment (straight-through):}
\begin{equation}
  c_v = \arg\max_b s_{v,b}
  \label{eq:gsb_hard}
\end{equation}

\textbf{Temperature annealing:}
\begin{equation}
  \tau^{(t+1)} = \max\!\bigl(\tau^{(t)} \cdot (1 - r_\text{anneal}),\; \tau_\text{min}\bigr)
  \quad (r_\text{anneal}=3\times10^{-4},\; \tau_\text{min}=0.5)
  \label{eq:gsb_anneal}
\end{equation}

\textbf{EMA usage tracking:}
\begin{equation}
  \hat{p}_b^{(t)} = 0.99 \cdot \hat{p}_b^{(t-1)} + 0.01 \cdot \frac{\#\{v: c_v=b\}}{|\text{batch}|}
  \label{eq:gsb_ema}
\end{equation}
\end{mathbox}

\subsection{Entropy Loss}

\begin{codebox}[title={codebook.py -- codebook\_entropy\_loss()}]
\begin{lstlisting}
def codebook_entropy_loss(self):
    p = self._ema_usage / (self._ema_usage.sum() + 1e-8)
    entropy = -(p * torch.log(p + 1e-8)).sum()
    max_entropy = math.log(self._num_embeddings)
    return 1.0 - (entropy / max_entropy)   # 0 = uniform, 1 = collapsed
\end{lstlisting}
\end{codebox}

\begin{mathbox}[title={Normalized Entropy Loss}]
\begin{equation}
  H(\hat{p}) = -\sum_{b=1}^B \hat{p}_b \log \hat{p}_b
  \in [0,\, \log B]
\end{equation}
\begin{equation}
  \mathcal{L}_\text{entropy} = 1 - \frac{H(\hat{p})}{\log B}
  \in [0, 1]
  \label{eq:entropy_loss}
\end{equation}
Giá trị 0 khi phân phối đều hoàn toàn ($\hat{p}_b = 1/B$ với mọi $b$),
giá trị 1 khi suy sụp hoàn toàn (một centroid nhận toàn bộ token).
\end{mathbox}

\begin{insightbox}
\textbf{Gradient flow:} Với \texttt{nn.Parameter}, cả \texttt{embedding\_k}
và \texttt{embedding\_v} nhận gradient từ attention output qua công thức:
\begin{equation}
  \frac{\partial \mathcal{L}}{\partial e_b} =
    \frac{\partial \mathcal{L}}{\partial h_\text{global}} \cdot
    \frac{\partial h_\text{global}}{\partial v_x} \cdot
    W_\text{value}^\top \cdot \mathbf{1}[c_v = b]
\end{equation}
Điều này khác VQ-EMA gốc: codebook học qua gradient descent
\emph{cộng thêm} EMA usage tracking cho entropy loss.
Hai cơ chế bổ sung cho nhau: gradient cập nhật hướng của codebook,
EMA tracking cập nhật trọng số sử dụng.
\end{insightbox}

% ======================================================================
\section{Module 5: GlobalModule với Centroid Attention}
% ======================================================================

\subsection{Kiến trúc và Công thức}

\begin{codebox}[title={model.py -- RelGTLayer.global\_forward() (rút gọn)}]
\begin{lstlisting}
def global_forward(self, x, batch_idx):
    q_x = self.lin_proj_g(x)             # [B, global_dim]
    k_x = self.vq.get_k()                # [B_centroids, global_dim]
    v_x = self.vq.get_v()                # [B_centroids, global_dim]
    
    q = self.lin_query_g(q_x)            # [B, h*d_h]
    k = self.lin_key_g(k_x)              # [B_c, h*d_h]
    v = self.lin_value_g(v_x)            # [B_c, h*d_h]
    
    q,k,v = rearrange to [h, n, d_h]
    dots = einsum("h i d, h j d -> h i j", q, k) * scale
    
    # Popularity bias
    dots = dots + log(centroid_count)    # [1, 1, B_c]
    attn = softmax(dots, dim=-1)
    out  = einsum("h i j, h j d -> h i d", attn, v)
    out  = rearrange to [B, h*d_h]
\end{lstlisting}
\end{codebox}

\begin{mathbox}[title={Công thức Global Attention với Popularity Bias}]
\textbf{Multi-head attention đến centroids:}
\begin{equation}
  \text{dots}_{h,i,b} =
    \frac{\langle W_q^h q_i,\; W_k^h e_b \rangle}{\sqrt{d_h}}
    + \underbrace{\log N_b}_{\text{popularity bias}}
  \label{eq:global_attn}
\end{equation}
với $N_b = \#\{v \in \text{toàn bộ dataset}: c_v = b\}$ (đếm tích lũy).

\textbf{Output:}
\begin{equation}
  h_\text{global}(v_i) = \bigoplus_{h=1}^H
    \sum_{b=1}^B \text{softmax}_b(\text{dots}_{h,i,:})\cdot W_v^h e_b
  \label{eq:global_out}
\end{equation}

\textbf{Cập nhật chỉ số centroid:}
\begin{equation}
  \mathbf{c}[\texttt{batch\_idx}] \leftarrow \arg\min_b \|q_v - e_b\|_2
  \quad (\text{được thực hiện bởi \texttt{vq.update(q\_x)}})
\end{equation}
\end{mathbox}

\begin{observation}
\textbf{c\_idx là global buffer $\in \mathbb{R}^{N_\text{nodes}}$:}
Buffer \texttt{c\_idx} được cập nhật mỗi training step:
\begin{center}
\texttt{self.c\_idx[batch\_idx] = x\_idx.squeeze().to(torch.long)}
\end{center}
Điều này tạo ra ``bộ nhớ dài hạn'' về việc node nào thuộc centroid nào.
Trong môi trường DDP với nhiều GPU, buffer này \textbf{không được đồng bộ}
qua \texttt{dist.broadcast} trong code hiện tại -- mỗi GPU có phiên bản
riêng của \texttt{c\_idx}, có thể dẫn đến inconsistency.
\end{observation}

% ======================================================================
\section{Module 6: CrossAttentionBridge (model.py)}
% ======================================================================

\subsection{Công thức Bidirectional Cross-Attention}

\begin{mathbox}[title={CrossAttentionBridge -- Công thức đầy đủ}]
Đầu vào: $h_l \in \mathbb{R}^{B \times d}$ (local), $h_g \in \mathbb{R}^{B \times d}$ (global).

\textbf{Local $\to$ Global attention:}
\begin{align}
  \alpha &= \sigma\!\left(\frac{(W_q^l h_l) \cdot (W_k^g h_g)}{\sqrt{d}}\right) \in (0,1)^B \\
  \tilde{h}_l &= h_l + \alpha \odot W_v^g h_g
  \label{eq:cab_l2g}
\end{align}

\textbf{Global $\to$ Local attention:}
\begin{align}
  \beta &= \sigma\!\left(\frac{(W_q^g h_g) \cdot (W_k^l h_l)}{\sqrt{d}}\right) \\
  \tilde{h}_g &= h_g + \beta \odot W_v^l h_l
  \label{eq:cab_g2l}
\end{align}

\textbf{Gated merge:}
\begin{align}
  \gamma &= \sigma\!\left(W_\text{gate}[\tilde{h}_l \| \tilde{h}_g]\right) \in (0,1)^{B \times d} \\
  m &= \gamma \odot \tilde{h}_l + (1-\gamma) \odot \tilde{h}_g
\end{align}

\textbf{Output với residual:}
\begin{equation}
  \boxed{h_\text{bridge} = \text{LN}\!\left(W_\text{out}\, \text{Dropout}(m) + h_l\right)}
  \label{eq:cab_out}
\end{equation}
\end{mathbox}

\begin{issuebox}
\textbf{Dot product đơn chiều (single-head):} Attention được tính bằng
dot product vô hướng $\langle q, k \rangle \in \mathbb{R}^B$ thay vì
multi-head attention. Điều này \emph{significantly} hạn chế khả năng
mô hình hóa tương tác phức tạp giữa local và global:
\begin{equation}
  \alpha = \sigma\!\left(\frac{1}{d}\sum_{i=1}^d q_i k_i\right)
  \quad \text{-- chỉ học một ``hướng tương tác''}
\end{equation}
So sánh với multi-head ($H$ heads):
\begin{equation}
  \alpha_h = \text{softmax}\!\left(\frac{Q_h K_h^\top}{\sqrt{d_h}}\right)
  \quad \text{-- học $H$ hướng tương tác độc lập}
\end{equation}
\end{issuebox}

% ======================================================================
\section{Hàm Mất mát và Auxiliary Losses}
% ======================================================================

\subsection{Tổng hợp Loss trong Training}

\begin{codebox}[title={main\_node\_ddp.py -- training loop (rút gọn)}]
\begin{lstlisting}
pred = model(...)
pred = pred.view(-1) if pred.size(1) == 1 else pred
loss = loss_fn(pred.float(), labels)    # ONLY task loss!
loss.backward()
\end{lstlisting}
\end{codebox}

\begin{issuebox}
\textbf{Auxiliary losses KHÔNG được sử dụng trong training!}

Mặc dù \texttt{model.get\_aux\_loss()} được định nghĩa và các
\texttt{\_aux\_losses} được thu thập sau mỗi forward pass (codebook entropy,
local-global agreement), chúng \textbf{không được cộng vào hàm mất mát}
trong \texttt{main\_node\_ddp.py}.

Tổng loss thực tế chỉ là:
\begin{equation}
  \mathcal{L}_\text{actual} = \mathcal{L}_\text{task}
  \quad \text{(chỉ task loss!)}
\end{equation}

Trong khi thiết kế dự kiến:
\begin{equation}
  \mathcal{L}_\text{intended} = \mathcal{L}_\text{task}
    + \lambda_1 \mathcal{L}_\text{entropy}
    + \lambda_2 \mathcal{L}_\text{local\text{-}global}
\end{equation}
\textbf{Đây là bug quan trọng} -- toàn bộ cơ chế entropy regularization
và local-global agreement không có tác dụng.
\end{issuebox}

\subsection{Local-Global Agreement Loss}

\begin{codebox}[title={model.py -- RelGTLayer.forward() -- full mode}]
\begin{lstlisting}
with torch.no_grad():
    cos_sim = F.cosine_similarity(out_local, out_global, dim=-1).mean()
self._aux_losses['local_global_agreement'] = (1.0 - cos_sim) * 0.01
\end{lstlisting}
\end{codebox}

\begin{mathbox}[title={Công thức}]
\begin{equation}
  \mathcal{L}_\text{agree} = 0.01 \times \left(1 - \frac{1}{B}\sum_{i=1}^B
    \frac{h_l^{(i)} \cdot h_g^{(i)}}{\|h_l^{(i)}\|\|h_g^{(i)}\|}\right)
\end{equation}
\end{mathbox}

\begin{insightbox}
\textbf{Dùng \texttt{torch.no\_grad()}}:  Loss này được tính với gradient bị
chặn từ phía cosine similarity (chỉ dùng để tracking/monitoring), không thực
sự ảnh hưởng đến tham số. Ngay cả nếu được cộng vào total loss, gradient
sẽ bằng 0 (do no\_grad wrapper). Đây có thể là lỗi thiết kế -- nếu muốn
alignment loss thực sự hoạt động, cần bỏ \texttt{no\_grad}.
\end{insightbox}

% ======================================================================
\section{Phân tích Hệ thống Lấy mẫu Dữ liệu (utils.py)}
% ======================================================================

\subsection{Chiến lược 2-hop với Temporal Constraint}

\begin{mathbox}[title={Bất đẳng thức thời gian trong lấy mẫu}]
\textbf{1-hop filter:}
\begin{equation}
  \mathcal{N}_1(v_i) = \{v_j \in \mathcal{N}_{1\text{-hop}}(v_i) : \tau(v_j) \le \tau(v_i)\}
\end{equation}

\textbf{2-hop filter:}
\begin{equation}
  \mathcal{N}_2(v_i) = \{v_k \in \mathcal{N}_{2\text{-hop}}(v_i)
    \setminus \mathcal{N}_1(v_i) : \tau(v_k) \le \tau(v_i)\}
\end{equation}

\textbf{Quyết định lấy mẫu:}
\begin{equation}
  \text{Neighbors} = \begin{cases}
    \text{Random sample}(\mathcal{N}_1 \cup \mathcal{N}_2,\; K-1)
      & \text{nếu } |\mathcal{N}| \ge K-1 \\
    \text{Random choice with replacement}(\mathcal{N},\; K-1)
      & \text{nếu } 0 < |\mathcal{N}| < K-1 \\
    \text{Global fallback sample}(\mathcal{V},\; K-1)
      & \text{nếu } |\mathcal{N}| = 0
  \end{cases}
\end{equation}
\end{mathbox}

\subsection{HDF5 Precomputation Pipeline}

\begin{table}[H]
\centering\small
\caption{Format lưu trữ HDF5 và kích thước dữ liệu}
\begin{tabular}{llll}
\toprule
\textbf{Dataset} & \textbf{Dtype} & \textbf{Shape} & \textbf{Nội dung} \\
\midrule
\texttt{types} & int16 & $[N, K]$ & Loại node cho mỗi token \\
\texttt{indices} & int32 & $[N, K]$ & Chỉ số local của node \\
\texttt{hops} & int8 & $[N, K]$ & Khoảng cách hop \\
\texttt{times} & float32 & $[N, K]$ & Thời gian tương đối (ngày) \\
\texttt{edges} & int16 & $[2, E_\text{total}]$ & Cạnh đồ thị con \\
\texttt{edges\_offsets} & uint64 & $[N+1]$ & Offset CSR cho edges \\
\bottomrule
\end{tabular}
\end{table}

\begin{issuebox}
\textbf{int16 cho edge index có thể gây tràn số:} \texttt{edges} được lưu
với dtype \texttt{int16}, có giá trị tối đa $2^{15}-1 = 32767$. Nhưng
$K = 300$ node/subgraph, vậy edge index từ 0 đến 299 -- không tràn số.
Tuy nhiên, nếu $K > 32767$ trong tương lai, đây sẽ là vấn đề.

Hơn nữa, trong \texttt{main\_node\_ddp.py} khi load:
\begin{center}
\texttt{batch["node\_indices"].to(device)}
\end{center}
Tensor int16 cần cast sang int64 trước khi dùng làm index. Code xử lý
điều này đúng với \texttt{.to(torch.int64)} ở đầu training loop.
\end{issuebox}

% ======================================================================
\section{Phân tích DDP và Training (main\_node\_ddp.py)}
% ======================================================================

\subsection{Gradient Clipping và Optimizer}

\begin{mathbox}[title={Cấu hình tối ưu hóa}]
\textbf{Optimizer:} Adam với learning rate scaling tuyến tính theo world size:
\begin{equation}
  \text{lr}_\text{eff} = \text{lr}_\text{base} \times |\mathcal{W}|,
  \quad \text{lr}_\text{base} = 10^{-4}
\end{equation}

\textbf{Gradient clipping:}
\begin{equation}
  \nabla \leftarrow \nabla \cdot \min\!\left(1,\; \frac{1}{\|\nabla\|_2}\right)
  \quad (\texttt{max\_norm} = 1.0)
\end{equation}

\textbf{Weight decay:} $\lambda = 10^{-5}$ (L2 regularization).
\end{mathbox}

\subsection{Evaluation Gathering trong DDP}

\begin{mathbox}[title={Phương pháp gather kết quả phân tán}]
Mỗi rank $r$ tính predictions trên phần dữ liệu của mình:
$\{(\text{idx}_i, \hat{y}_i)\}_{i \in S_r}$.

Rank 0 gather tất cả:
\begin{equation}
  \text{all\_preds} = \bigoplus_{r=0}^{|\mathcal{W}|-1}
    \{(\text{idx}_i, \hat{y}_i)\}_{i \in S_r}
\end{equation}
\begin{equation}
  \text{final\_preds}[\text{idx}] = \hat{y}_{\text{idx}}
  \quad \forall \text{idx} \in [0, N_\text{split})
\end{equation}
\end{mathbox}

% ======================================================================
\section{Tổng hợp: Ma trận Điểm mạnh / Điểm yếu}
% ======================================================================

\begin{table}[H]
\centering\small
\renewcommand{\arraystretch}{1.3}
\caption{Đánh giá toàn diện từng module của RelGT++}
\begin{tabularx}{\textwidth}{lXXc}
\toprule
\textbf{Module} & \textbf{Điểm mạnh} & \textbf{Điểm yếu / Vấn đề} & \textbf{Mức độ} \\
\midrule
TypeEncoder & Đơn giản, hiệu quả, tương đương toán học với paper & Không có & ✓✓✓ \\
HopEncoder & Dịch chuyển $+1$ xử lý đúng padding & Chỉ 3 giá trị hop, không học được & ✓✓ \\
TimeEncoder & Xử lý token không có thời gian qua mask vector & PE dùng sin/cos theo giây -- chu kỳ có thể không phù hợp & ✓✓ \\
TfsEncoder & Grouping-encode-scatter rất hiệu quả ($10^4\times$ speedup) & NaN clip tới $10^6$ có thể gây gradient bùng nổ & ✓✓✓ \\
GNNPEEncoder & Stochastic init tăng expressiveness & Variance cao inference, GIN dim $d/4$ & ✓✓ \\
CrossModalFusion & Gate per-modality thích nghi với context & Dead modality problem, $4$M tham số & ✓✓✓ \\
EncoderLayer & Flash attn, DropPath đúng, xavier init & DropPath toàn batch không phải per-sample & ✓✓✓ \\
GatedFFN (SwiGLU) & Gradient flow tốt, expansion $4\times$ & BN trên token dim -- ít ổn định hơn LN & ✓✓✓ \\
MultiViewReadout & Ba góc nhìn bổ sung, residual với seed & -- & ✓✓✓ \\
GumbelSoftCodebook & Gumbel exploration, temperature annealing, gradient flow & Temperature giảm không dựa trên metric thực, EMA decoupled từ gradient & ✓✓✓ \\
CrossAttnBridge & Bidirectional, gated merge có lý thuyết & Attention đơn chiều (dot-product vô hướng), no\_grad trong agreement loss & ✓✓ \\
GlobalAttn & Popularity bias tốt cho cân bằng centroid & c\_idx không sync trong DDP & ✓✓ \\
Training loop & Gradient clip, SyncBN, gather distributed eval & Auxiliary losses không được dùng (bug!) & ✓ \\
\bottomrule
\end{tabularx}
\end{table}

\subsection{Top 3 Vấn đề Ưu tiên Sửa}

\begin{enumerate}
  \item \textbf{[Quan trọng nhất]} Auxiliary losses không được cộng vào loss:
        Sửa bằng cách thêm \texttt{loss = loss + model.module.get\_aux\_loss()}
        sau dòng \texttt{loss = loss\_fn(...)} trong training loop.
  \item \textbf{[Quan trọng]} Agreement loss tính với \texttt{no\_grad}:
        Bỏ wrapper \texttt{with torch.no\_grad()} trong
        \texttt{RelGTLayer.forward()} khi tính cosine similarity.
  \item \textbf{[Trung bình]} c\_idx không sync qua DDP:
        Thêm \texttt{dist.all\_reduce(self.c\_idx, op=dist.ReduceOp.SUM)}
        (hoặc dùng \texttt{all\_gather}) sau mỗi batch update trong global forward.
\end{enumerate}

% ======================================================================
\section{Kết luận}
% ======================================================================

RelGT++ là một triển khai kỹ thuật công phu với nhiều cải tiến thực sự
so với RelGT gốc. Đặc biệt, sự kết hợp của GumbelSoftCodebook (gradient
flow qua codebook), GatedFFN (SwiGLU activation), CrossModalGatedFusion
(per-modality adaptive weighting), và MultiViewReadout (3 aggregation views)
tạo ra một kiến trúc biểu cảm cao.

Tuy nhiên, phân tích từ mã nguồn phát hiện \textbf{ba vấn đề kỹ thuật
quan trọng}: auxiliary losses không được kích hoạt, agreement loss bị
chặn gradient, và codebook index không đồng bộ trong môi trường phân tán.
Sửa ba vấn đề này (ước tính $\sim$5 dòng code) có thể cải thiện đáng kể
hiệu suất huấn luyện mà không cần thay đổi kiến trúc.

\begin{thebibliography}{9}
\bibitem{relgt25}
Dwivedi et al. (2025). \textit{Relational Graph Transformer}. arXiv:2505.10960.
\bibitem{swiglu20}
Shazeer, N. (2020). \textit{GLU Variants Improve Transformer}. arXiv:2002.05202.
\bibitem{stodepth20}
Huang, G. et al. (2016). \textit{Deep Networks with Stochastic Depth}. ECCV.
\bibitem{gumbel17}
Jang, E. et al. (2017). \textit{Categorical Reparameterization with Gumbel-Softmax}. ICLR.
\bibitem{gin19}
Xu, K. et al. (2019). \textit{How Powerful are Graph Neural Networks?}. ICLR.
\bibitem{pytorchframe}
Hu, W. et al. (2024). \textit{PyTorch Frame}. arXiv:2404.00776.
\end{thebibliography}

\end{document}
